\chapter{Preliminaries}
\label{chapter:preliminaries}

This chapter develops the necessary preliminaries for what follows by recalling the notions of principal bundles, classifying spaces, and characteristic classes. We assume that the reader is familiar with the concept of a fiber bundle with structure group. A classical reference for this is, for example, \cite{steenrod1951topology}. Moreover, we presuppose some basic familiarity with homotopy and homology theory. The first section, on principal bundles, draws mainly on \cite{davis2001lecture} and \cite{Mitchell2011PrincipalBundles}. In the second section, on classifying spaces, we follow \cite{ebert_bordism_script}. Finally, the third section, on characteristic classes, is based primarily on \cite{Hatcher2003VectorBundles}, with additional material drawn from \cite{LawsonMichelsohn1989SpinGeometry} and \cite{Frenck2026Doubling}.

\section{Principal Bundles and Associated Bundles}
\label{subsec:principal_bundles}

\begin{definition}
    \label{bundles:def:principal_bundle}
    Let $G$ be a topological group. A $G$-principal bundle is a fiber bundle with fiber $G$ and structure group $G$ acting by left translations.
\end{definition}

The total space of a $G$-principal bundle $P \xrightarrow{\pi} B$ carries a natural right action of $G$ defined as follows. Around $y \in P$, choose a local trivialisation
\[
\phi \colon U \times G \to \pi^{-1}(U)
\]
and let $\phi^{-1}(y) = (\pi(y),h)$. Then set 
\[
y \cdot g := \phi(\pi(y),hg).
\]
This is well defined since, by definition, $G$ acts on itself by \textit{left} translations under a change of local trivialisation. 
A morphism of principal bundles should respect this right action, thus we define:

\begin{definition}
    \label{bundles:def:principal_bundles_morphism}
    A morphism from a $G$-principal bundle $P \to B$ to an $H$-principal bundle $P^\prime \to B^\prime$ is a triple $(F,f, \rho)$ of continuous maps $F \colon P \to P^\prime$ and $f \colon B \to B^\prime$ together with a homomorphism $\rho \colon G \to H$, such that 
    \[\begin{tikzcd}
	P & {P^\prime} \\
	B & {B^\prime}
	\arrow["F", from=1-1, to=1-2]
	\arrow[from=1-1, to=2-1]
	\arrow[from=1-2, to=2-2]
	\arrow["f", from=2-1, to=2-2]
    \end{tikzcd}\]
    commutes and such that $F$ is $\rho$-equivariant, that is $F(x \cdot g) = F(x) \cdot \rho(g)$ for all $x \in P$ and $g \in G$. If $G=H$, we assume $\rho$ to be the identity, and if $B = B^\prime$, we assume $f$ to be the identity.
\end{definition}

The condition of equivariance in the above definition is actually quite strong, as illustrated by the following proposition.

\begin{proposition}
    \label{bundles:prop:morphisms_of_principal_bundles}
    Any morphism of $G$-principal bundles over the same base space is an isomorphism.
\end{proposition}

\begin{proof}
    First suppose we are given a morphism 
    \[  
        F \colon B \times G \to B \times G
    \]
    of trivial bundles. Then by definition $F$ is of the form 
    \[
    F(x,g) = (x,\phi_x(g))
    \]
    for a continuous family of maps $\phi_x \colon G \to G$. By equivariance, $\phi_x(g) = \phi_x(e) g$, and thus an inverse of $F$ is given by
    \[
    F^{-1}(x,g) := (x, \phi_x(e)^{-1} g).
    \]
    Since every $G$-principal bundle is locally trivial, this shows that we can always find local inverses in the general case. As inverses are unique, they fit together to define a global inverse.
\end{proof}

To any fiber bundle with structure group $G$, one can canonically associate a $G$-principal bundle by a general construction that allows one to change the fiber of a given bundle. We now recall this construction of changing the fiber. 
Let $p \colon E \to B$ be a fiber bundle with fiber $F$ and structure group $G$ and let $F^\prime$ be another space upon which $G$ acts effectively. Then we can construct a new bundle as follows. Choose an open cover $(U_\alpha)_{\alpha \in A}$ of $B$ together with local trivialisations $\phi_\alpha \colon U_\alpha \times F \to p^{-1}(U_\alpha)$. Let $\tau_{\alpha\beta} \colon U_\alpha \cap U_\beta \to G$ be the corresponding transition functions defined by 
\[
(\phi_\alpha^{-1} \circ \phi_\beta)(x,f) = (x, \tau_{\alpha\beta}(x) \cdot f).
\]
Define 
\[
E^\prime := \left( \coprod_{\alpha \in A} U_\alpha \times F^\prime \right)_{\raisebox{5pt}{\scalebox{1.5}{$/$}}\raisebox{5pt}{\scalebox{1.05}{$\sim$}}}
\]
where the equivalence relation $\sim$ is defined by
\[
(b,f^\prime,\beta) \sim (b,\tau_{\alpha\beta}(b) \cdot f^\prime,\alpha)
\]
and let $p^\prime \colon E^\prime \to B, \: [(b,f^\prime,\alpha)] \mapsto b$. Then the maps
\[
    \psi_\alpha \colon U_\alpha \times F^\prime \to (p^\prime)^{-1}(U_\alpha), \; (b,f^\prime) \mapsto [(b,f^\prime,\alpha)].
\]
define local trivialisations of $p^\prime$ with the same transition functions $\tau_{\alpha\beta}$, giving $p^\prime \colon E^\prime \to B$ the structure of a fiber bundle with fiber $F^\prime$ and structure group $G$. We say that this new bundle is obtained from $E \to B$ by \textit{changing the fiber} from $F$ to $F^\prime$ (with the $G$-action on $F^\prime$ being implicit). Note that we can also obtain $E \to B$ from $E^\prime \to B$ by changing the fiber from $F^\prime$ back to $F$. Thus no information is lost in this process. Considering the special case $F^\prime = G$ with $G$ acting on itself by left translations, we call the bundle obtained from $E \to B$ by changing the fiber from $F$ to $G$ the \textit{associated principal bundle}. Morally speaking, if we want to understand the bundle $E \to B$, it is sufficient to understand its associated principal bundle, since this contains the same information, as we have seen above. The benefit of studying principal bundles is that they are highly symmetric objects, and thus usually easier to understand. For example, many constructions simplify for principal bundles, such as the concept of universal bundles, as we will see in the next section. Another such instance is the construction for changing the fiber from above. In the case of principal bundles this can be carried out more directly and has the benefit of also allowing for non-effective group actions. This is known as the \textit{Borel Construction}.

\begin{lemma}
    \label{bundles:lem:borel_construction}
    Let $p \colon P \to B$ be a $G$-principal bundle and $F$ a space upon which $G$ acts from the left. Then 
    \[
    E := \modulo{(P \times F)}{\sim}, \quad (y,f) \sim (y\cdot g, g^{-1}f)
    \]
    together with the map $\pi \colon E \to B$ given by $[(y,f)] \mapsto p(y)$ is a fiber bundle with fiber $F$ and structure group $G$. Moreover, if the action of $G$ on $F$ is effective, this bundle is isomorphic to the bundle obtained from $P \to B$ by changing the fiber from $G$ to $F$ using the construction outlined above.
\end{lemma}

\begin{proof}
    To see that $P \times_G F$ is locally trivial, let $\psi \colon U \times G \to P$ be a local trivialisation. Then the homeomorphism 
    \[
    \psi \times \mathrm{id}_F \colon U \times G \times F \to p^{-1}(U) \times F
    \]
    descends to a homeomorphism
    \[
    \tilde{\kappa} \colon \modulo{(U \times G \times F)}{\sim} \to \modulo{(p^{-1}(U) \times F)}{\sim} = \pi^{-1}(U),
    \]
    where the equivalence relation on $U \times G \times F$ is given by $(x,g,f) \sim (x,gh,h^{-1}f)$. Moreover, the map 
    \[
    U \times G \times F \to U \times F, \; (x,g,f) \mapsto (x,gf)
    \]
    descends to a homeomorphism $\modulo{(U \times G \times F)}{\sim} \to U \times F$ with inverse given by $(x,f) \mapsto [(x,e,f)]$. Under this homeomorphism, $\tilde{\kappa}$ corresponds to a homeomorphism $\kappa \colon U \times F \to \pi^{-1}(U)$, which is explicitly given by $(x,f) \mapsto [(\psi(x,e),f)]$. Since
    \[\begin{tikzcd}
	{U \times F} & {\pi^{-1}(U)} \\
	{} & U
	\arrow["{\kappa}", from=1-1, to=1-2]
	\arrow["{\mathrm{proj.}}"', from=1-1, to=2-2]
    \arrow["\pi", from=1-2, to=2-2]
    \end{tikzcd}\]
    commutes, this shows that $\pi \colon E \to B$ is locally trivial. To see that, for effective $G$-actions, this construction agrees with the change-of-fiber construction introduced above, we compute the transition functions of the local trivialisations just defined. Let $U_\alpha, U_\beta \subseteq B$ be two trivialising neighbourhoods for $P \to B$ with local trivialisations $\psi_\alpha, \psi_\beta$ and transition function $\tau_{\alpha\beta} \colon U_\alpha \cap U_\beta \to G$. Moreover, let $\kappa_\alpha, \kappa_\beta$ be the corresponding local trivialisations for $E \to B$ as constructed above. Then we compute 
    \begin{align*}
        (\kappa_\alpha^{-1} \circ \kappa_\beta)(x,f) & = \kappa_\alpha^{-1}([(\psi_\beta(x,e),f)]) = (x,\tau_{\alpha\beta}(x) f).
    \end{align*}
    Thus the statement follows from the fact that fiber bundles, where the structure group acts effectively, are uniquely determined by the transition functions of local trivialisations.
\end{proof}

We denote the bundle $E \to B$ obtained from the above construction by $P \times_G F \to B$. The Borel construction is natural with respect to pullbacks, as stated in the following lemma.

\begin{lemma}
        \label{bundles:lem:natruality_of_changing_fiber}
        Let $p \colon P \to B$ be a $G$-principal bundle, $F$ a space upon which $G$ acts from the left and $f \colon B^\prime \to B$ a continuous map. Then 
        \[
        (f^\ast P) \times_G F \cong f^\ast(P \times_G F).
        \]
\end{lemma}

\begin{proof}
    An explicit model for the pullback bundle is 
    \[
    f^\ast P := \{(y,b^\prime) \in P \times B^\prime \mid p(y) = f(b^\prime)\}.
    \]
    Then 
    \[
    (f^\ast P) \times_G F = \{(y,b^\prime,f) \in P \times B^\prime \times F \mid p(y) = f(b^\prime)\} / \sim,
    \]
    for $(y, b^\prime, f) \sim (y \cdot g, b^\prime, g^{-1} \cdot f)$. Moreover, 
    \begin{align*}
    f^\ast(P \times_G F) & = \{([y,f],b^\prime) \in P \times_G F \times B^\prime \mid p(y) = f(b^\prime)\} \\
    & \cong \{(y,f,b^\prime) \in P \times F \times B^\prime \mid p(y) = f(b^\prime)\} / \sim
    \end{align*}
    for $(y,f,b^\prime) \sim (y \cdot g, g^{-1} \cdot f, b^\prime)$. Clearly these are isomorphic bundles.
\end{proof}

\begin{remark}
    \label{bundles:rem:naturality_change_of_fiber}
    We remark that the action of $G$ on $F$ in the above lemma need not be effective. If we restrict ourselves to effective actions, a more general naturality result holds for the change-of-fiber construction of arbitrary bundles. Namely, first pulling a bundle back under some continuous map and then changing the fiber amounts to the same as first changing the fiber and then pulling back. Comparing the transition functions of the local trivialisations in both cases shows that they in fact agree.
\end{remark}

Let's look at some concrete examples for the theory developed above.

\begin{example}
    \label{bundles:ex:associated_bundles}
    \begin{enumerate}[label=(\roman*)]
        \item Let $\pi \colon E \to B$ be a real rank $n$ vector bundle over a paracompact base space $B$. Choosing a bundle metric gives $E \to B$ the structure of a fiber bundle with fiber $\R^n$ and structure group $O(n)$. The associated principal bundle is explicitly given by the \textit{orthonormal frame bundle $\mathrm{O}(E)$ of $E$}, whose fiber over a point $x \in B$ is the space of all ordered orthonormal bases of $E_x$. The topology on $O(E)$ is the subspace topology $O(E) \subseteq E \times \ldots \times E$. Local trivialisations of the canonical projection map $p \colon \mathrm{O}(E) \to B$ can be defined as follows. Choose an orthonormal local frame $e_1,\ldots,e_n$ of $E$ defined over an open set $U \subseteq B$. Then a trivialisation $\phi \colon U \times \mathrm{O}(n) \to p^{-1}(U)$ is given by
        \[
            (x,A) \mapsto \left(\sum_{i=1}^n a_{i1}e_i(x),\ldots,\sum_{i=1}^n a_{in}e_i(x)\right),
        \]
        where $A = (a_{ij})$. So $\phi(x,A)$ is the basis of $E_x = \pi^{-1}(x)$, whose coordinates in the frame $(e_i)$ are precisely the columns of the matrix $A$. The transition functions of these trivialisations agree with the transition functions of the trivialisations of $E$. Since the transition functions uniquely determine the bundle, we conclude that $O(E)$ is in fact isomorphic to the associated principal bundle as defined at the beginning of this section. Analogously, if $E \to B$ is a complex rank $n$ vector bundle, choosing a hermitian metric on $E$ allows us to reduce the structure group to $U(n)$, and the corresponding associated principal bundle is the \textit{unitary frame bundle}. 
        % \item An orientation of a real vector bundle allows us to further reduce the structure group to $SO(n)$. Similarly, the associated principal bundle is then given by the \textit{oriented orthonormal frame bundle $\so(E)$ of $E$}, whose fiber over a point $x \in B$ is the space of all oriented ordered orthonormal bases of $E_x$. 
        \item \label{bundles:ex:associated_bundles_cpn} Let $p \colon \tau_1^n \to \CP^n$ be the tautological complex line bundle equipped with the canonical hermitian metric. Since this bundle is of rank $1$, its unitary frame bundle is simply the subspace $U(\tau_1^n) \subseteq \tau_1^n$ of all vectors of norm $1$, together with the restriction of $p$ to this subspace. Recall that
        \[
        \tau_1^n = \{(L,z) \in \CP^n \times \C^{n+1} \mid z \in L\}.
        \]
        So, explicitly,
        \[
        U(\tau_1^n) = \{(L,z) \in \tau_1^n \mid \abs{z} = 1\}.
        \]
        Consider the diffeomorphism 
        \[
        U(\tau_1^n) \to S^{2n+1}, \; (L,z) \mapsto z
        \]
        with inverse $S^{2n+1} \to U(\tau_1^n),\; z \mapsto (\pi(z),z)$, where $\pi \colon S^{2n+1} \to \CP^n$ is the canonical projection. Under this diffeomorphism the restriction of $p$ to $U(\tau_1^n)$ corresponds to $\pi$. Hence $\pi$ defines a $U(1) \cong S^1$-principal bundle, which agrees with the associated principal bundle of $\tau_1^n \to \CP^n$.
        \item \label{bundles:ex:associated_bundles_part3} Let $E \to B$ be an $S^1$-principal bundle over a smooth manifold $B$ without boundary. We can change the fiber of $E$ from $S^1$ to $D^2$ to obtain another bundle $DE \to B$, where $S^1$ acts on $D^2$ by rotations. By definition of the change-of-fiber construction, trivialisations $U \times D^2 \to DE$ have the same transition functions as the corresponding trivialisations of $E \to B$. Moreover, they restrict to trivialisations $U \times S^1 \to \partial DE$ of the boundary. Since the action of $S^1$ on $D^2$ restricts to the action of $S^1$ on itself, we see that $E \cong \partial DE$. So, in particular, the total space of any $S^1$-principal bundle is null-bordant. We will make use of this fact in \hyperref[chapter:main_construction]{Chapter 4}.
    \end{enumerate}
\end{example}


\section{Universal Bundles and Classifying Spaces}
A natural question is which fiber bundles with specified fiber and structure group can exist over a given base space. As before, this problem turns out to be easier for principal bundles. We assume all base spaces to be CW-complexes in this section.

\begin{definition}
    A $G$-principal bundle $E \to B$ is called \textit{universal} if $E$ is weakly contractible as a topological space.
\end{definition}

The significance of this definition is established by the following theorem. We say that two morphisms of $G$-principal bundles are $G$-homotopic if they are homotopic through morphisms of $G$-principal bundles.

\begin{theorem}
    \label{bundles:thm:universal_principal_bundle}
    Let $E \to B$ be a universal $G$-principal bundle. Then for any other $G$-principal bundle $P \to X$ there exists a morphism $P \to E$ of $G$-principal bundles, which is unique up to $G$-homotopy.
\end{theorem}

Before showing the theorem, let us discuss its implications. Firstly, we can equivalently formulate the theorem by saying that $E \to B$ is a terminal object in the homotopy category of $G$-principal bundles over CW-complexes and thus, in particular, uniquely determined (up to $G$-homotopy equivalence). In view of this, one usually writes $EG := E$ and $BG := B$, and refers to $EG \to BG$ as \textit{the} universal $G$-principal bundle. The base space $BG$ is called the \textit{classifying space} of $G$. Since any bundle morphism $(F,f) \colon P \to EG$ induces a morphism $P \to f^\ast EG$ (and vice versa, which is true for arbitrary fiber bundles), we conclude from Proposition \ref{bundles:prop:morphisms_of_principal_bundles} that 
\[
P \cong f^\ast EG.
\]
We call $f$ a \textit{classifying map} for $P$. By homotopy invariance of the pullback, we get a well-defined bijection 
\[
\phi_X \colon [X,BG] \to \mathcal{P}_G(X), \; [f] \mapsto [f^\ast EG],
\]
where $\mathcal{P}_G(X)$ denotes the set of isomorphism classes of $G$-principal bundles over $X$. The existence result in Theorem \ref{bundles:thm:universal_principal_bundle} shows that $\phi_X$ is surjective and the uniqueness result shows that it is injective. Moreover, $\phi_X$ is clearly natural in $X$. Thus we have shown:

\begin{corollary}
\label{bundles:cor:universal_bundle_defines_natural_isomorphism}
    The maps $\phi_X$ define a natural isomorphism of contravariant functors 
    \[  
        [-,BG], \mathcal{P}_G(-) \colon \mathsf{HCW} \to \mathsf{Set},
    \]
    where $\mathsf{HCW}$ denotes the homotopy category of CW-complexes.
\end{corollary}

Theorem \ref{bundles:thm:universal_principal_bundle} will be a simple application of the following theorem.

\begin{theorem}
    \label{bundles:thm:extension_theorem}
    Let $P \to X$ be a $G$-principal bundle and $A \subseteq X$ a subcomplex. Then any morphism $P|_A \to EG$ of $G$-principal bundles extends to a morphism $P \to EG$ of $G$-principal bundles. 
\end{theorem}

The proof of this theorem is based on an inductive, cell-by-cell approach. We refer the reader to \cite[Theorem 2.7.1]{ebert_bordism_script}. 

\begin{proof}[Proof of Theorem \ref{bundles:thm:universal_principal_bundle}]
    Applying Theorem \ref{bundles:thm:extension_theorem} to $A = \emptyset$ shows the existence of a morphism $P \to EG$. Now suppose we are given morphisms $F_0,F_1 \colon P \to EG$. Consider the $G$-principal bundle 
    \[
    \pi \times \mathrm{id} \colon P \times [0,1] \to X \times [0,1],
    \]
    where $\pi \colon P \to X$ is the bundle projection (this is just the pullback bundle of $P$ under the projection $X \times [0,1] \to X$). Then $F_0$ and $F_1$ define a morphism of $G$-principal bundles 
    \[
    (F_0,F_1) \colon P \times \{0,1\} = (P \times [0,1])|_{X \times \{0,1\}}\to EG,
    \]
    which by Theorem \ref{bundles:thm:extension_theorem} can be extended to $P \times [0,1]$, thus yielding a $G$-homotopy between $F_0$ and $F_1$.
\end{proof}

Of course, for this theory to be of any use, we need such universal bundles to exist. The following general existence theorem was first shown by Milnor \cite{milnor1956universalII}.

\begin{theorem}
    For any topological group $G$ there exists a universal $G$-principal bundle $EG \to BG$, where $BG$ is a CW-complex.
\end{theorem}

A priori, the space $BG$ constructed in \cite{milnor1956universalII} may not be a CW-complex. However, if we choose a CW-approximation $f \colon X \to BG$, the total space of the bundle $f^\ast EG \to X$ is still weakly contractible. This follows by applying the five lemma to the following diagram.
\[\begin{tikzcd}
	{\pi_{k+1}(X)} & {\pi_k(G)} & {\pi_k(f^\ast EG)} & {\pi_k(X)} & {\pi_{k-1}(G)} \\
	{\pi_{k+1}(BG)} & {\pi_k(G)} & {\pi_k(EG)} & {\pi_k(BG)} & {\pi_{k-1}(G)}
	\arrow[from=1-1, to=1-2]
	\arrow["\cong", from=1-1, to=2-1]
	\arrow[from=1-2, to=1-3]
	\arrow["{\mathrm{id}}", from=1-2, to=2-2]
	\arrow[from=1-3, to=1-4]
	\arrow[from=1-3, to=2-3]
	\arrow[from=1-4, to=1-5]
	\arrow["\cong", from=1-4, to=2-4]
	\arrow["{\mathrm{id}}", from=1-5, to=2-5]
	\arrow[from=2-1, to=2-2]
	\arrow[from=2-2, to=2-3]
	\arrow[from=2-3, to=2-4]
	\arrow[from=2-4, to=2-5]
\end{tikzcd}\]

We can now use these classification results for principal bundles to obtain analogous results for general fiber bundles by appealing to the change-of-fiber construction discussed in the preceding section.

\begin{corollary}
    \label{bundles:cor:universal_fiber_bundles}
    Let $G$ be a topological group acting effectively on a space $F$. Then for any CW-complex $X$ the map 
    \[
    [X,BG] \to \mathcal{F}_{F,G}(X), \; \; [f] \mapsto \left[f^\ast(EG \times_G F)\right]
    \]
    is a natural bijection, where $\mathcal{F}_{F,G}(X)$ denotes the set of isomorphism classes of fiber bundles with fiber $F$ and structure group $G$ over $X$.
\end{corollary}

\begin{proof}
    By Corollary \ref{bundles:cor:universal_bundle_defines_natural_isomorphism}, the map 
    \[
    [X,BG] \to \mathcal{P}_G(X), \; \; [f] \mapsto [f^\ast EG]
    \]
    is a natural bijection, and by Lemma \ref{bundles:lem:natruality_of_changing_fiber}
    \[
    \mathcal{P}_G(X) \to \mathcal{F}_{F,G}(X), \; \; [P] \mapsto \left[P \times_G F\right]
    \]
    is a natural bijection. Composing both thus yields the natural bijection 
    \[
    [X,BG] \to \mathcal{F}_{F,G}(X), \; \; [f] \mapsto \left[(f^\ast EG) \times_G F\right].
    \]
    So the statement follows by again applying Lemma \ref{bundles:lem:natruality_of_changing_fiber}.
\end{proof}

Again let us discuss a few examples.

\begin{example}
    \label{bundles:ex:examples_for_classifiying_spaces}
    \begin{enumerate}[label=(\roman*)]
        \item[]
        \item The universal cover $\R \to S^1$ is a $\Z$-prinicpal bundle, $\Z$ acting on itself by addition. Since $\R$ is contractible we obtain $B\Z \cong S^1$ and $E\Z \cong \R$.
        \item \label{bundles:ex:homotopy_groups_of_BG} From the long exact sequence of homotopy groups 
        \[\begin{tikzcd}
        {\pi_k(EG)} & {\pi_k(BG)} & {\pi_{k-1}(G)} & {\pi_{k-1}(EG)}
        \arrow[from=1-1, to=1-2]
        \arrow[from=1-2, to=1-3]
        \arrow[from=1-3, to=1-4]
        \end{tikzcd}\]
        we conclude that $\pi_k(BG) \cong \pi_{k-1}(G)$ for $k \geq 1$. Moreover, $BG$ is connected since it is the image of the connected space $EG$ under $EG \to BG$.
        \item The universal $O(k)$-prinicpal bundle $EO(k) \to BO(k)$ also has an explicit description. Consider the Stiefel manifold $V_k^n \subseteq \R^{n\times k}$, which is defined as the compact space of orthonormal $k$-tuples of vectors in $\R^n$. Moreover let $\mathrm{Gr}_k^n$ be the Grassmanian of $k$-planes in $\R^n$, i.e. the set of $k$-dimensional linear subspaces of $\R^n$. The topology on $\mathrm{Gr}_k^n$ is induced by the projection map
        \[
        p \colon V_k^n \to \mathrm{Gr}_k^n, \; (v_1,\ldots,v_k) \mapsto \mathrm{span}(v_1,\ldots,v_k)
        \] 
        Identifying a linear subspace $V \in \mathrm{Gr}_k^n$ with $\R^k$ yields an identification of the fiber $p^{-1}(V)$ with $O(k)$. Moreover the identifications $V \cong \R^k$ can be chosen consistently in a small neighbourhood around $V$, providing local trivialisations for $p$. The transition of two such trivialisations is realized by actions of $O(k)$, hence $p$ defines an $O(k)$-principal bundle.
        Now the canonical inclusion $\R^n \to \R^{n+1}$ induces inclusions $V_k^n \to V_k^{n+1}$ and $\mathrm{Gr}_k^n \to \mathrm{Gr}_k^{n+1}$. Consider the colimits over these inclusions 
        \[
        V_k^\infty := \operatorname*{colim}_n\,(V_k^n), \quad \mathrm{Gr}_k^\infty := \operatorname*{colim}_n\,(\mathrm{Gr}_k^n)
        \] 
        and the induced map $V_k^\infty \to \mathrm{Gr}_k^\infty$. One can choose trivialisations of $V_k^n \to \mathrm{Gr}_k^n$ for varying $n$, which fit together to induce a trivialisation in the colimit, such that $V_k^\infty \to \mathrm{Gr}_k^\infty$ defines an $O(k)$-principal bundle. We claim that $V_k^\infty$ is weakly contractible. Consider the map 
        \[
        V_k^n \to S^{n-1}, \; (v_1, \ldots, v_k) \mapsto v_1.
        \]
        Similiar arguments as above shows that this map is locally trivial with fiber $V_{k-1}^{n-1}$. Hence we can consider the long exact sequence of homotopy groups 
        \[\begin{tikzcd}[column sep=scriptsize]
        {\pi_{i+1}(S^{n-1})} & {\pi_i(V_{k-1}^{n-1})} & {\pi_i(V_n^k)} & {\pi_i(S^{n-1})}
        \arrow[from=1-1, to=1-2]
        \arrow[from=1-2, to=1-3]
        \arrow[from=1-3, to=1-4]
        \end{tikzcd}\]
        So $\pi_i(V_k^n) \cong \pi_i(V_{k-1}^{n-1})$ for $i \leq n-3$. Applying this inductively yields 
        \[
        \pi_i(V_k^n) \cong \pi_i(V_1^{n-(k-1)}) = \pi_i(S^{n-k+1}) = 0
        \] 
        for $i \leq n-k-1$. Now from compactness of $S^i$ and the fact that $V_k^n$ is closed in $V_k^{n+1}$ we conclude that the image of any map $S^i \to V_k^\infty$ is contained in some finite stage $V_k^n$. Therefore $\pi_i(V_k^\infty) = 0$ for all $i \geq 0$, such that the universal $O(k)$-principal bundle is explicitly given by $V_k^\infty \to \mathrm{Gr}_k^\infty$.

        Similiarly we get an explicit description of the universal $SO(k)$-prinicpal bundle. Let $\tilde{\mathrm{Gr}}_k^n$ be the set of tuples consisting of a $k$-dimensional linear subspace of $\R^n$ together with a chosen orientation. The topology on $\tilde{\mathrm{Gr}}_k^n$ is induced by 
        \[
        p \colon V_k^n \to \tilde{\mathrm{Gr}}_k^n, \; (v_1,\ldots,v_k) \mapsto (\mathrm{span}(v_1,\ldots,v_k), \mathrm{o}_{(v_1,\ldots,v_k)}),
        \] 
        where $\mathrm{o}_{(v_1,\ldots,v_k)}$ is the unique orientation on $\mathrm{span}(v_1,\ldots,v_k)$, such that $(v_1,\ldots,v_k)$ is positively oriented. Similiar to above we see that $p$ defines an $SO(k)$-principal bundle. Taking the colimit, we obtain the universal $SO(k)$-principal bundle $V_k^\infty \to \tilde{\mathrm{Gr}}_k^\infty$.

        \item \label{bundles:ex:examples_for_classifiying_spaces_universal_vector_bundle} Any real $k$-dimensional vector bundle $E \to X$ can be considered as a fiber bundle with fiber $\R^k$ and structure group $O(k)$. So Corollary \ref{bundles:cor:universal_fiber_bundles} and the example above show that every such $E$ is the pullback of the universal real $k$-dimensional vector bundle $V_k^\infty \times_{O(k)} \R^k$ under some map $X \to \mathrm{Gr}_k^\infty$. We show that this universal vector bundle also has a more explicit description, namely as the tautological bundle $\gamma_k \to \mathrm{Gr}_k^\infty$. Recall that this bundle is defined as 
        \[
        \gamma_k := \{(X,v) \in \mathrm{Gr}_k^\infty \times \R^\infty \mid v \in X\}
        \] 
        together with the canonical projection map to $\mathrm{Gr}_k^\infty$. An element in $V_k^\infty$ is a $k$-tuple of orthonormal vectors in some $\R^n$, which we will subsequently regard as the columns of an $n \times k$ matrix. Then the map 
        \[
        V_k^\infty \times \R^k \to \gamma_k, \;\; (A,v) \mapsto (\mathrm{im}(A), Av)
        \]
        descends to the quotient, defining a vector bundle map $V_k^\infty \times_{O(k)} \R^k \to \gamma_k$. This map is clearly a linear isomorphism on fibers, hence an isomorphism of vector bundles. 

        Similiarly, oriented real $k$-dimensional vector bundles are classified by maps to $\tilde{\mathrm{Gr}}_k^\infty$. We claim that the universal oriented real $k$-dimensional vector bundle is given by the tautological bundle $\gamma_k^{SO} \to \tilde{\mathrm{Gr}}_k^\infty$ defined by 
        \[
        \gamma_k^{SO} := \{(X,\mathrm{o}_X,v) \in \tilde{\mathrm{Gr}}_k^\infty \times \R^k \mid v \in X\}
        \]
        together with the canonical projection map to $\tilde{\mathrm{Gr}}_k^\infty$. To this end consider the map 
        \[
        V_k^\infty \times \R^k \to \gamma_k^{SO}, \;\; (A,v) \mapsto (\mathrm{im}(A), \mathrm{o}_A, Av),
        \]
        where $\mathrm{o}_A$ is the unique orientation on $\mathrm{im}(A)$, such that $A \colon \R^k \to \mathrm{im}(A)$ is orientation preserving. Again this map descends to the quotient defining an isomorphism $V_k^\infty \times_{SO(k)} \R^k \to \gamma_k^{SO}$.



         
        
        
        % $BO(n)$ and $EO(n)$ explicit via Stiefel Manifolds (compare n lab) + $EO(n) \times_{O(n)} \R^n \cong \gamma_n$ for $\gamma_n$ tautological bundle over grassmanain + orientable vector bundles classfied by BSO(n) (universal vector bundle over this denoted by $\gamma_n^{\so}$)
        \item \label{bundles:ex:examples_for_classifiying_spaces_universal_complex_vector_bundle}Completely analogous to the real case we can also define the complex Stiefel manifolds $V_k^\infty(\C)$ and complex Grassmanians $\mathrm{Gr}_k^\infty(\C)$ and show that the universal $U(k)$-prinicpal bundle is given by $V_k^\infty(\C) \to \mathrm{Gr}_k^\infty(\C)$. The universal complex $k$-dimensional vector bundle is the tautological bundle $\tau_k \to \mathrm{Gr}_k^\infty(\C)$. 
        \item \label{bundles:es:examples_for_classifying_spaces_special_case}As special cases of the preceeding examples we obtain
        \[
        B\Z_2 \cong \RP^\infty, \quad BS^1 \cong \CP^\infty. 
        \]
    \end{enumerate}
\end{example}



\section{Stiefel-Whitney and Chern-Classes}

In the preceding chapter we discussed the question of how bundles can be classified. Another fundamental question is, of course, how to determine whether two given bundles are isomorphic. As usual in topology, a partial solution to this problem is provided by defining computable invariants. For vector bundles the most important invariants are so-called \textit{characteristic classes}. The general definition of a characteristic class is as follows. Let $\K = \R, \C$ and write $\mathrm{Vect}_\K(B)$ for the set of isomorphism classes of $\K$-vector bundles over a space $B$. Then $\mathrm{Vect}_\K(-)$ defines a contravariant functor $\mathsf{Top} \to \mathsf{Set}$. A characteristic class $c$ for $\K$-vector bundles with values in a cohomology theory $H^\ast$ is a natural transformation from the functor $\mathrm{Vect}_\K(-)$ to the functor $H^\ast$, regarded as a functor with values in $\mathsf{Set}$. In more concrete terms, $c$ assigns to each $\K$-vector bundle $E \to B$ a cohomology class $c(E) \in H^\ast(B)$, depending only on the isomorphism class of $E$, such that for any continuous map $f \colon B' \to B$ the following naturality condition holds:
\[
c(f^\ast E) = f^\ast c(E).
\]
Using the theory of classifying spaces developed above, it follows that for real (resp. complex) $k$-dimensional vector bundles, the characteristic class $c$ is uniquely determined by its value on the universal bundle $\gamma_k$ (resp. $\tau_k$) discussed in Example~\ref{bundles:ex:examples_for_classifiying_spaces}~\ref{bundles:ex:examples_for_classifiying_spaces_universal_vector_bundle} (resp.~\ref{bundles:ex:examples_for_classifiying_spaces_universal_complex_vector_bundle6}).

For real vector bundles the most important characteristic classes are the \textit{Stiefel-Whitney classes} and for complex vector bundles the \textit{Chern classes}. Both are characterized by a few axioms.

\begin{definition}[Stiefel-Whitney Classes, HATCHER]
    \label{bundles:def:Steifel_Whitney_Classes}
    There is a unique sequence of characteristic classes $w_1, w_2, \ldots$ for real vector bundles with values in $H^\ast(-;\Z_2)$, such that 
    \begin{enumerate}[label=(\roman*)]
        \item $w_i(E) \in H^i(B;\Z_2)$ for any real bundle $E \to B$,
        \item \label{bundles:eq:prop2_for_stiefel_whitney_classes} $w_i(E) = 0$ for $i > \mathrm{rank}(E)$,
        \item \label{bundles:eq:prop3_for_stiefel_whitney_classes} $w(E_1 \oplus E_2) = w(E_1) \smile w(E_2)$, where $w = 1 + w_1 + w_2 + \ldots$,
        \item \label{bundles:eq:prop4_for_stiefel_whitney_classes} For the universal real line bundle $\gamma_1 \to \RP^\infty$, $w_1(\gamma_1)$ is the generator of $H^1(\RP^\infty;\Z_2) \cong \Z_2$.
    \end{enumerate}
\end{definition}

Note that by property \ref{bundles:eq:prop2_for_stiefel_whitney_classes} the sum $w = 1 + w_1 + w_2 + \ldots$ is always finite. The class $w(E)$ is called the \textit{total Stiefel-Whitney class of $E$}. We sometimes write $w_0(E) := 1 \in H^0(B;\Z_2)$ to simplify notation.

\begin{definition}[Chern Classes, HATCHER]
    There is a unique sequence of characteristic classes $c_1, c_2, \ldots$ for complex vector bundles with values in $H^\ast(-;\Z)$, such that 
    \begin{enumerate}[label=(\roman*)]
        \item $c_i(E) \in H^{2i}(B;\Z)$ for any complex bundle $E \to B$,
        \item $c_i(E) = 0$ for $i > \mathrm{rank}(E)$,
        \item $c(E_1 \oplus E_2) = c(E_1) \smile c(E_2)$, where $c = 1 + c_1 + c_2 + \ldots$,
        \item For the universal complex line bundle $\tau_1 \to \CP^\infty$, $c_1(\tau_1)$ is a generator of $H^2(\CP^\infty;\Z) \cong \Z$ specified in advance.
    \end{enumerate}
\end{definition}

Again the \textit{total Chern class} $c = 1 + c_1 + c_2 + \ldots$ is well-defined and we write $c_0(E) := 1 \in H^0(B;\Z)$. Note that in both definitions, property \ref{bundles:eq:prop4_for_stiefel_whitney_classes} uniquely determines the values of the respective characteristic classes on line bundles and, together with property \ref{bundles:eq:prop3_for_stiefel_whitney_classes}, prescribes their values on bundles that split as sums of line bundles. The values on arbitrary bundles can then be deduced from a general splitting principle [??]. We refer the reader to [HATCHER] for a detailed proof of existence and uniqueness.

Not only are the definitions of Stiefel-Whitney and Chern classes formally quite similiar, but there is also a direct connection between the two.

\begin{proposition}[Hatcher Proposition 3.8]
    \label{bundles:prop:connection_chern_stiefel_whitney_classes}
    Regarding a complex rank $n$ vector bundle $E \to B$ as a real rank $2n$ vector bundle, $w_{2i +1}(E) = 0$ and $w_{2i}(E)$ is the image of $c_i(E)$ under the coefficient homomorphism $H^{2i}(B;\Z) \to H^{2i}(B;\Z_2)$.  
\end{proposition}

Let us consider some examples.

\begin{example}
    \label{bundles:ex:characteristic_classes}
    \begin{enumerate}[label=(\roman*)]
        \item[]
        \item Let $\iota \colon \RP^n \to \RP^\infty$ be the canonical inclusion. Then the pullback $\iota^\ast \gamma_1$ is isomorphic to the tautological line bundle $\gamma_1^n \to \RP^n$. Since $\iota$ induces an isomorphism on cohomology we conclude from naturality that $w_1(\gamma_1^n) = \iota^\ast w_1(\gamma_1)$ is the generator of $H^1(\RP^n;\Z_2) \cong \Z_2$. Similiarly, $c_1(\tau_1^n)$ of the tautological complex line bundle $\tau_1^n \to \CP^n$ is a generator of $H^2(\CP^n;\Z) \cong \Z$.
        \item Let $\epsilon \to B$ be a trivial real bundle over some space $B$. Then $\epsilon \cong f^\ast \epsilon$ for a constant map $f \colon B \to B$ and thus $w_i(\epsilon_n) = f^\ast w_i(\epsilon) = 0$ for all $i \geq 1$. Similiarly all Chern classes for complex trivial bundles vanish.
        \item \label{bundles:ex:characteristic_classes_part3} Let $TS^n \to S^n$ be the tangent bundle of the sphere and let $\nu_n \to S^n$ be the normal bundle with respect to the standard embedding $S^n \to \R^{n+1}$. Since $\nu_n$ and $TS^n \oplus \nu_n$ are trivial we conclude from the example above and property \ref{bundles:eq:prop3_for_stiefel_whitney_classes} in Definition \ref{bundles:def:Steifel_Whitney_Classes} that 
        \[
        1 = w(TS^n \oplus \nu_n) = w(TS^n) \smile w(\nu_n) = w(TS^n) \smile 1 = w(TS^n).
        \]
        Thus $w_i(TS^n) = 0$ for all $i \geq 1$.
    \end{enumerate}
\end{example}

In the last example we saw that although all Stiefel-Whitney classes of a bundle vanish, it need not be trivial. However for one-dimensional real (resp. complex) bundles it turns out that the first Stiefel-Whitney (resp. Chern) class already completely determines the bundle:

\begin{theorem}[Hathcer, Proposition 3.10]
    \label{bundles:thm:line_bundles_determined_by_chern_class}
    Let $X$ be a CW-complex. Then the maps 
    $$w_1 \colon \mathrm{Vect}_\R(X) \to H^1(X;\Z_2)$$ 
    and 
    $$c_1 \colon \mathrm{Vect}_\C(X) \to H^2(X;\Z)$$ 
    are bijections. 
\end{theorem}

% This allows us to define line bundles by only prescribing their first Stiefel-Whitney class (resp. Chern class), which is often technically easier thank directly writing down this bundle. 
While the first Stiefel-Whitney class of higher dimensional bundles does not identify their isomorphism type it does characterize another important invariant, namely orientability:

\begin{proposition}[Hatcher, Prop 3.11]
    \label{bundles:prop:orientable_iff_w1_vanishes}
    Let $X$ be a CW-complex. A real vector bundle $E \to X$ is orientable if and only if $w_1(E) = 0$.
\end{proposition}

Similiarly vanishing of the second Stiefel Whitney class corresponds to the existence of another structure, namely that of a \textit{spin structure}. 

\begin{definition}
    \begin{enumerate}[label=(\roman*)]
        \item For $n \geq 2$ the \textit{spin group} $\mathrm{Spin}(n)$ is defined as the total space of the non-trivial two-sheeted cover $\rho \colon \mathrm{Spin}(n) \to \mathrm{SO}(n)$. If $n \geq 3$ this is the universal cover and in the case $n=2$ we have $\spin(2) \cong \so(2)$. For $n=1$ we define $\spin(1) := \Z_2$, the total space of the two-sheeted cover $\rho \colon \Z_2 \to \so(1)$.
        \item Let $E \to B$ be an oriented real vector bundle equipped with a fiber metric and let $P_{\mathrm{SO}(n)}$ be its oriented orthonormal frame bundle. A \textit{spin structure} on $E$ is a $\mathrm{Spin}(n)$-principal bundle $P_{\mathrm{Spin}(n)}$ together with a map 
        \[
        \phi \colon P_{\mathrm{Spin}(n)} \to P_{\mathrm{SO}(n)}
        \]
        of principal bundles, which is $\rho$-equivariant (compare Definition \ref{bundles:def:principal_bundles_morphism}).
        \item An orientable manifold $M$ is called \textit{spin} if its tangent bundle $TM$ admits a spin structure and is called \textit{totally non-spin} if the universal cover of $M$ is not spin.
    \end{enumerate}
\end{definition}

\begin{remark}
    When $n=1$, $P_{\so(1)} \cong X$ and a spin structure is simply a two-sheeted cover of $X$.
\end{remark}
Note that the existence of a spin structure does not depend on the chose orientation or fiber metric on $E$, since different choices yield isomorphic oriented orthonormal frame bundles. In fact the only obstruction to the existence of a spin structure is the second Stiefel-Whitney class, as mentioned above.

\begin{theorem}[Lawson,Micheaslon] %in der quelle steht es glaub ich nur für X eine maniigfaltigkeit, suche noch andere quelle
    \label{bundles:thm:spin_iff_w2_vanishes}
    Let $X$ be a CW-complex. An orientable real vector bundle $E \to X$ admits a spin structure if and only if $w_2(E) = 0$. 
\end{theorem}

\begin{example}
    \begin{enumerate}[label=(\roman*)]
        \item All spheres $S^n$, $n\geq 1$ are spin by Example \ref{bundles:ex:characteristic_classes} \ref{bundles:ex:characteristic_classes_part3}.
        \item \label{bundles:ex:cp2_not_spin} The standard example for non-spin manifolds are all even dimensional complex projective spaces $\CP^{2k}$, see [MilnorStasheff??].
    \end{enumerate}
\end{example}

\begin{remark}
    \label{bundles:rem:classifying_map_for_spin_bundles}
    Let $E \to X$ be a real vector bundle over a CW-complex. If $E$ is orientable, we have seen that this bundle is classified by a map $X \to B\mathrm{SO}(n)$. Similiarly, if a spin structure exists, $E$ is classified by a map $X \to B\mathrm{Spin}(n)$, which can be seen as follows. Let $\phi \colon P_{\mathrm{Spin}(n)} \to P_{\mathrm{SO}(n)}$ be a spin structure on $E$ and let $f \colon X \to B\mathrm{Spin}(n)$ be a classifiying map for $P_{\mathrm{Spin}(n)}$. Define a vector bundle 
    \[
    \gamma_n^\spin := E\mathrm{Spin}(n) \times_{\mathrm{Spin}(n)} \R^n
    \]
    using Lemma \ref{bundles:lem:borel_construction}, where $\mathrm{Spin}(n)$ acts on $\R^n$ by virtue of the covering map $\rho \colon \mathrm{Spin}(n) \to \mathrm{SO}(n)$ and the canonical action of $\so(n)$. From Lemma \ref{bundles:lem:natruality_of_changing_fiber} we obtain
    \[
    f^\ast \gamma_n^\spin \cong P_{\spin(n)} \times_{\spin(n)} \R^n.
    \]
    Moreover, the spin structure $\phi$ induces an isomorphism $P_{\spin(n)} \times_{\spin(n)} \R^n \to P_{\so(n)} \times_{\so(n)} \R^n \cong E$ of vector bundles by
    \[
    [(x,v)] \mapsto [(\phi(x),v)].
    \]
    This map is clearly a linear bijection on fibers, hence an isomorphism of vector bundles. 
\end{remark}

Recall that a manifold $M$ is orientable if and only if its tangent bundle $TM$ restricted to any embedded loop $S^1 \subseteq M$ is trivial. Or in other words ``$TM$ does not contain any Moebius Bands''. A similiar characterization can be given for spin manifolds, which is why the property of being spin is sometimes referred to as a higher dimensional analogue of orientability.

\begin{proposition}[Georg Paper, Proposition 2.6]
    \label{bundles:prop:spin_characterisation}
    Let $M$ be a connected oriented manifold of dimension at least $5$.
    \begin{enumerate}[label=(\roman*)]
        \item \label{bundles:eq:characterization_spin} $M$ is spin if and only if for every embedded surface $S \subseteq M$ the restriction $TM|_S$ of the tangent bundle to $S$ is trivial.
        \item \label{bundles:eq:characterization_totally_non_spin} $M$ is totally non-spin if and only if there exists an embedding $ \iota \colon S^2 \to M$, such that $w_2(\iota^\ast TM) \neq 0$. In particular, the restriction $\iota^\ast TM$ of the tangent bundle to this embedded sphere is non-trivial.
    \end{enumerate}
\end{proposition}

\begin{proof}
    Since we will not be making use of the first characterization, we shall only proof part \ref{bundles:eq:characterization_totally_non_spin} of the proposition and refer the reader to [Georg Paper] for a proof of part \ref{bundles:eq:characterization_spin}. Let $\tilde{M}$ be the universal cover of $M$. Since $T\tilde{M}$ does not admit a spin structure we conclude from Theorem \ref{bundles:thm:spin_iff_w2_vanishes} that $w_2(T\tilde{M}) \neq 0$. So by the universal coefficient theorem there exists a homology class $a \in H_2(\tilde{M};\Z_2)$ with 
    \[
    \langle w_2(T\tilde{M}), a\rangle \neq 0.
    \]
    Once more applying the universal coefficient theorem as well as the Hurewicz theorem we obtain an isomorphism
    \[
    \pi_2(\tilde{M}) \otimes \Z_2 \to H_2(\tilde{M};\Z) \otimes \Z_2 \cong H_2(\tilde{M};\Z_2). 
    \]
    Explicitly, for some class $\alpha \in \pi_2(\tilde{M})$ represented by a map $f \colon S^2 \to \tilde{M}$ the above isomorphism maps $\alpha \otimes 1$ to $f_\ast [S^2]_{\Z_2}$, where $[S^2]_{\Z_2}$ is the $\Z_2$-fundamental class of $S^2$. Therefore we can write $a = f_\ast [S^2]_{\Z_2}$ for some $f \colon S^2 \to \tilde{M}$. Let $p \colon \tilde{M} \to M$ be the projection and define $\iota := p \circ f \colon S^2 \to M$. Since $\mathrm{dim}(M) \geq 5$ we may assume $\iota$ to be an embedding. We compute 
    \[
    \langle w_2(\iota^\ast TM), [S^2]_{\Z_2}\rangle = \langle f^\ast w_2(p^\ast TM), [S^2]_{\Z_2}\rangle = \langle w_2(T\tilde{M}), f_\ast [S^2]_{\Z_2}\rangle = \langle w_2(T\tilde{M}), a \rangle \neq 0,
    \]
    where we used $p^\ast TM \cong T \tilde{M}$ and naturality of the Kronecker pairing in the second equation. Therefore $w_2(\iota^\ast TM) \neq 0$.
\end{proof}

Finally we record another scenario in which the vanishing of all Stiefel-Whitney classes of a bundle is sufficient for triviality.

\begin{lemma}
    \label{bundles:lem:criterion_for_triviality}
    Let $E \to S$ be a real vector bundle of rank at least $3$ over a two dimensional CW-complex $S$. Then $E$ is trivial if and only if $w_1(E) = w_2(E) = 0$.
\end{lemma}

\begin{proof}
    Combining Theorem \ref{bundles:thm:spin_iff_w2_vanishes} and Remark \ref{bundles:rem:classifying_map_for_spin_bundles} we conclude that there exists a classifiying map $f \colon S \to B\spin(n)$ for $E$, where $n$ is the rank of $E$. From Example \ref{bundles:ex:examples_for_classifiying_spaces} \ref{bundles:ex:homotopy_groups_of_BG} it follows that $\pi_k(B\spin(n)) = 0$ for $k \geq 2$. Since $S$ is two dimensional the map $f$ is homotopic to a constant ([Hatcher, Lemma 4.6]), therefore $E \cong f^\ast \xi_\spin$ is trivial.
\end{proof}